% =====================================================================
%  07  支配方程式（ヘッダー＝支配方程式）
%  source_memo: 04_reaction.tex §4.2.4 支配方程式（物質・エネルギー・Ergun）/
%               §4.2.5 スイング設計（N_swing,sets・V_cat・N_parallel・V_vessel）/
%               §4.2 量論行列(eq:stoich_matrix)・成分ベクトル
%  方針：見出しは小チップ。速度式は §反応工程で既出のため載せない。
%        各列に「式（上）＋記号表（下）」を縦積みし、中央を全高の縦線で仕切る。
%        左＝モデル（式＋記号表：ベクトル/行列 F・r・ν・ΔH・Cp を収録）。
%        右＝スイング設計（式＋記号表）。表は 1記号=1行・2群・群間に縦罫線。
%        モデル式は短いので左表を直下から始め余白を回収し表を最大化。色なし・丸括弧禁止。
% =====================================================================
\begin{frame}[t]

  \vspace{0.15em}

  \begin{columns}[T,onlytextwidth]
    % ===== 左：モデル（式＋記号表）================================
    \begin{column}{0.49\textwidth}
      \fbox{\scriptsize モデル}\par
      \vspace{0.75em}
      {\fontsize{5.5}{6.6}\selectfont [物質収支]}\par
      {\centering$\displaystyle\frac{d\boldsymbol{F}}{dz}=(1-\varepsilon)\,A_{\mathrm{cross}}\,\boldsymbol{\nu}\,\boldsymbol{r}$\par}
      \vspace{0.95em}
      {\fontsize{5.5}{6.6}\selectfont [エネルギー収支・断熱]}\par
      {\centering$\displaystyle\frac{dT}{dz}=-\frac{(1-\varepsilon)\,A_{\mathrm{cross}}\,\boldsymbol{\Delta H}^{\!\top}\boldsymbol{r}}{\boldsymbol{F}^{\!\top}\boldsymbol{C}_p}$\par}
      \vspace{0.95em}
      {\fontsize{5.5}{6.6}\selectfont [Ergun]}\par
      {\centering\resizebox{\linewidth}{!}{$\displaystyle\frac{dP}{dz}=-\Big[\,150\frac{(1-\varepsilon_b)^2\mu u}{\varepsilon_b^3(\phi d_p)^2}
           +1.75\frac{(1-\varepsilon_b)\rho u^2}{\varepsilon_b^3\phi d_p}\,\Big]$}\par}
      \vspace{0.6em}
      {\color{black!25}\rule{\linewidth}{0.4pt}}\par
      \vspace{0.45em}
      % --- 左表：モデルの記号（ベクトル/行列＋支配式スカラー）2群 ---
      \resizebox{\linewidth}{!}{%
       \renewcommand{\arraystretch}{1.25}%
       \begin{tabular}{@{}l l l @{\hspace{1em}}|@{\hspace{1em}} l l l@{}}
         \toprule
         $\boldsymbol{F}$        & モル流量 & mol/s
           & $\varepsilon_b$ & 粒子間空隙率 & -- \\
         $\boldsymbol{r}$        & 反応速度 & mol/m$^3$s
           & $d_p$          & 粒径          & m \\
         $\boldsymbol{\nu}$      & 量論行列 & --
           & $\phi$          & 形状係数      & -- \\
         $\boldsymbol{\Delta H}$ & 反応熱   & kJ/mol
           & $T$            & 温度          & K \\
         $\boldsymbol{C}_p$      & 熱容量   & J/molK
           & $P$            & 全圧          & Pa \\
         $z$                     & 軸位置   & m
           & $u$            & 空塔速度      & m/s \\
         $A_{\mathrm{cross}}$    & 断面積   & m$^2$
           & $\rho$          & 密度          & kg/m$^3$ \\
         $\varepsilon$           & 床/容器体積比 & --
           & $\mu$          & 粘度          & Pa$\cdot$s \\
         \bottomrule
       \end{tabular}}
    \end{column}

    % ===== 中央：全高の仕切り線 ==================================
    \begin{column}{0.02\textwidth}
      \centering
      {\color{black!35}\rule[-6.9cm]{0.6pt}{7.05cm}}
    \end{column}

    % ===== 右：スイング設計（式＋記号表）==========================
    \begin{column}{0.49\textwidth}
      \fbox{\scriptsize スイング設計}\par
      \vspace{0.45em}
      {\fontsize{5.5}{6.6}\selectfont [再生追従のセット数]}\par
      {\centering$\displaystyle N_{\mathrm{swing,sets}}=\left\lceil\frac{t_{\mathrm{regen}}}{t_{\mathrm{cyc}}}\right\rceil+1$\par}
      \vspace{0.35em}
      {\fontsize{5.5}{6.6}\selectfont [触媒体積]}\par
      {\centering$\displaystyle V_{\mathrm{cat}}=(1-\varepsilon)\,z_{\mathrm{cat}}\,\frac{\pi D^2}{4}$\par}
      \vspace{0.35em}
      {\fontsize{5.5}{6.6}\selectfont [容器分割の並列基数]}\par
      {\centering$\displaystyle N_{\mathrm{parallel}}=\max\!\left(\left\lceil\frac{V_{\mathrm{cat}}}{V_{\mathrm{cat,max}}}\right\rceil,\,1\right)$\par}
      \vspace{0.35em}
      {\fontsize{5.5}{6.6}\selectfont [容器1基の体積　CAPEX 基準]}\par
      {\centering$\displaystyle V_{\mathrm{vessel}}=\frac{z_{\mathrm{cat}}}{N_{\mathrm{parallel}}}\cdot\frac{\pi D^2}{4}$\par}
      \vspace{0.5em}
      {\color{black!25}\rule{\linewidth}{0.4pt}}\par
      \vspace{0.4em}
      % --- 右表：スイング設計の記号 2群 ---
      \resizebox{\linewidth}{!}{%
       \renewcommand{\arraystretch}{1.25}%
       \begin{tabular}{@{}l l l @{\hspace{1em}}|@{\hspace{1em}} l l l@{}}
         \toprule
         $z_{\mathrm{cat}}$   & 触媒層長 & m
           & $V_{\mathrm{cat}}$    & 触媒体積       & m$^3$ \\
         $D$                  & 内径     & m
           & $V_{\mathrm{vessel}}$ & 容器体積       & m$^3$ \\
         $t_{\mathrm{cyc}}$   & 反応時間 & s
           & $V_{\mathrm{cat,max}}$ & 最大触媒体積  & m$^3$ \\
         $t_{\mathrm{regen}}$ & 再生時間 & s
           & $N_{\mathrm{parallel}}$ & 並列基数     & -- \\
         & & & $N_{\mathrm{swing,sets}}$ & スイングセット数 & -- \\
         \bottomrule
       \end{tabular}}
    \end{column}
  \end{columns}

\end{frame}
